% Generated draft: requires mathematical and source review.
\documentclass[11pt]{article}
\usepackage[margin=25mm]{geometry}
\usepackage{fontspec}
\setmainfont{Latin Modern Roman}
\usepackage{amsmath,amssymb,mathtools,unicode-math}
\setmathfont{Latin Modern Math}
\usepackage{enumitem,needspace,xurl,hyperref}
\setlength{\parindent}{0pt}
\setlength{\parskip}{5pt}
\setlength{\emergencystretch}{3em}
\newcommand{\N}{\mathbb{N}}
\newcommand{\Z}{\mathbb{Z}}
\newcommand{\Q}{\mathbb{Q}}
\newcommand{\R}{\mathbb{R}}
\newcommand{\C}{\mathbb{C}}
\newcommand{\F}{\mathbb{F}}
\newcommand{\A}{\mathbb{A}}
\newcommand{\PP}{\mathbb{P}}
\newcommand{\E}{\mathbb{E}}
\newcommand{\eps}{\varepsilon}
\newcommand{\dd}{\,\mathrm{d}}
\begin{document}
\section*{UnsolvedMath ID: 700572}
\textbf{Source problem code:} catalogue entry
\subsection{Definitions and mathematical statement}
Work in ZFC, including the Axiom of Choice; the Continuum Hypothesis is not assumed. Let \(\Omega\) be a countable set, and let \(\mathfrak m\) be a cardinal satisfying
\[
\aleph_0\leq\mathfrak m\leq 2^{\aleph_0}.
\]
Write \(\operatorname{Sym}(\Omega)\) for the group of all bijections of \(\Omega\), and \(\mathfrak P(\Omega)\) for its power set. A permutation group \(G\leq\operatorname{Sym}(\Omega)\) acts on \(\mathfrak P(\Omega)\) by
\[
g\cdot A=\{g(a):a\in A\}.
\]
Thus \(A,B\subseteq\Omega\) belong to the same orbit if and only if \(B=g\cdot A\) for some \(g\in G\).

The source asks whether, for every cardinal \(\mathfrak m\) in the stated interval, there exists such a group \(G\) with
\[
\left|\mathfrak P(\Omega)/G\right|=\mathfrak m.
\]
Here \(\mathfrak P(\Omega)/G\) denotes the set of orbits. The source says only ``countable,'' without explicitly saying ``countably infinite.'' The prescribed infinite number of orbits requires \(\Omega\) to be infinite, so the intended domain appears to be \(|\Omega|=\aleph_0\). This restriction is an unresolved wording qualification, rather than an explicitly supplied hypothesis.

The supplied comment dated 2001 reports that the answer is affirmative under Martin's Axiom, citing S. Shelah and S. Thomas (1988), and describes the question as still open in ZFC. These are supplied historical assessments, not independently verified conclusions.
\subsection{Short English statement}
For a countably infinite set, can every cardinal between \(\aleph_0\) and \(2^{\aleph_0}\), inclusive, occur as the number of orbits of a permutation group acting on its subsets? The source's use of ``countable'' needs clarification. Its 2001 comment reports a positive answer under Martin's Axiom and an open status in ZFC.
\subsection{Sources}
\begin{itemize}
\item Local UnsolvedMath registry: supplied complete record, permanent ID \(700572\), exact source code \texttt{catalogue entry}, titled \emph{Kourovka Notebook Problem 9.39}. The registry labels the problem open; this label is supplied source data, not independent verification.
\item P. M. Neumann, problem 9.39, \emph{The Kourovka Notebook (Unsolved Problems in Group Theory)}, 9th issue (1984). The supplied background identifies the 21st issue edition (2026), \texttt{arXiv:1401.0300v48}; the supplied collection description names E. I. Khukhro and V. D. Mazurov as editors.
\item S. Shelah and S. Thomas, \emph{Bull. London Math. Soc.} \textbf{20}, no. 4 (1988), 313--318, cited in the supplied comment of 2001 for the affirmative answer under Martin's Axiom. No article title is supplied. The comment's literature assessment is historical source material and has not been independently verified.
\end{itemize}
\end{document}
